\subsection{An Artin transfer from the Eisenstein field}
\label{sec:artin-induced-cutoff}
The following is a corollary of the fixed-field Hecke theorem.
Artin--Hecke reciprocity, induction of $L$-functions, and the
conductor formula are classical; no additional analytic estimate
for a nonabelian representation is claimed.

Let $K=\Q(\sqrt{-3})$, let $G_F=\operatorname{Gal}(\overline{\Q}/F)$,
and let $\sigma$ be the nontrivial automorphism of $K$.
For a continuous finite-order character $\chi$ of $G_K$, also
write $\chi$ for its primitive Hecke character, using arithmetic
Frobenius so that $\chi(\operatorname{Frob}_{\mathfrak p})
=\chi(\mathfrak p)$ at an unramified prime.
This correspondence and its local-factor compatibility are stated in
\citet[Theorems~4.2--4.3 and Proposition~4.4, pp.~52--53;
Proposition~5.2, p.~56]{horawaLfunctions}.
Extend ideal values by zero at the conductor primes.
Choose complex conjugation $j$ lifting $\sigma$, put
$\chi^\sigma(g)=\chi(jgj^{-1})$, and suppose
\[
 \chi\ne\chi^\sigma,\qquad
 \rho=\operatorname{Ind}_{G_K}^{G_\Q}\chi,\qquad
 q_\rho=3N\mathfrak f_\chi,\qquad c_\rho=2\pi/\sqrt{q_\rho}.
\]
In particular $\chi$ is nonprincipal.

On $L^2((0,\infty),du)$ use the Hankel transform
$\mathcal H_{c_\rho}$ defined above and, for fixed $0<\nu<1$, set
\[
 \Pi_{\rho,\Lambda}^{(\nu)}
  =\mathbf1_{[\nu,1]}
   (P_\Lambda\mathcal H_{c_\rho}P_\Lambda
                     \mathcal H_{c_\rho}P_\Lambda),
 \qquad P_\Lambda=\mathbf1_{(0,\Lambda)}.
\]
Summing over nonzero integral ideals of $K$, define the norm-line
theta map and its full Gram image projection by
\[
 E_\rho\psi(t)=e^{t/2}\sum_{\mathfrak a\ne0}
                     \chi(\mathfrak a)\psi(N\mathfrak a\,e^t),
 \qquad
 Q_{\rho,\Lambda}^{(\nu)}
     =P_{E_\rho\operatorname{Ran}\Pi_{\rho,\Lambda}^{(\nu)}}.
\]
There is no principal subtraction or constraint in this family.
For the nontrivial zeros $\omega=\beta_\omega+i\gamma_\omega$
of $L_\Q(s,\rho)$, with multiplicities $m_\omega$, put
\[
 \mathfrak b_\rho
   =\sum_\omega m_\omega\mu_{\beta_\omega-1/2,\gamma_\omega},
 \quad
 \mathfrak B_\rho(h)=\int H(i\tau)\,d\mathfrak b_\rho(\tau),
 \quad
 \mathfrak W_\rho(h)=\sum_\omega m_\omega H(\omega-1/2).
\]
Here $H(z)=\int h(t)e^{zt}dt$ and $\mu_{a,\gamma}$ is the
probability measure of Section~\ref{sec:harmonic}: its density is
$|a|/(\pi((\tau-\gamma)^2+a^2))$ when $a\ne0$, and it is
$\delta_\gamma$ when $a=0$.

\begin{corollary}[The cutoff for quadratic-induced Artin functions]
\label{cor:artin-induced-cutoff}
The representation $\rho$ has finite image, is irreducible of
dimension two, has Artin conductor $q_\rho$, and satisfies
\begin{equation}\label{eq:artin-induced-L}
                   L_\Q(s,\rho)=L_K(s,\chi)
\end{equation}
with every finite Euler factor, including the ramified factors.
This function is entire. Its real-place gamma factor is
$\Gamma_\R(s)\Gamma_\R(s+1)=\Gamma_\C(s)$, where
$\Gamma_\R(s)=\pi^{-s/2}\Gamma(s/2)$ and
$\Gamma_\C(s)=2(2\pi)^{-s}\Gamma(s)$.

Replacing $\chi$ by $\chi^\sigma$ leaves the raw selection,
theta map, and projection $Q_{\rho,\Lambda}^{(\nu)}$ unchanged.
For every fixed $\rho$ in this family, fixed $0<\nu<1$, and
$h\in C_c^\infty(\R)$,
\begin{equation}\label{eq:artin-induced-trace}
 \lim_{\Lambda\to\infty}
 \left(2\log\Lambda\,h(0)
      -\Tr(Q_{\rho,\Lambda}^{(\nu)}V(h))\right)
                         =\mathfrak B_\rho(h).
\end{equation}
The threshold $\nu=1/2$ is included, and
$\operatorname{rank}\Pi_{\rho,\Lambda}^{(\nu)}
\sim2c_\rho\Lambda/\pi$.

Put $L=\log\Lambda$, $S_L=\mathbf1_{[-L,L]}$, and let
$(z_{i,\Lambda})$ be an orthonormal basis of the theta image.
The nonnegative measures
\begin{equation}\label{eq:artin-induced-measures}
 d\mu_{\rho,\Lambda}(\tau)=
 \left(2L-\sum_i|\widehat{S_Lz_{i,\Lambda}}(\tau)|^2\right)
                                      \frac{d\tau}{2\pi}
\end{equation}
converge to $\mathfrak b_\rho$ on every Schwartz test.
For some $C,A$ depending on $\rho,\nu$, their masses on
$[T-1,T+1]$ are at most $C(1+|T|)^A$, uniformly for all
sufficiently large $\Lambda$ and all real $T$.
No zero-location, simplicity, or separation hypothesis is used.

If $L_\Q(s,\rho)$ is zero-free for $\Re s>\theta$,
$1/2\leq\theta<1$, put $\delta=\theta-1/2$. Then
\begin{equation}\label{eq:artin-induced-band}
 |\mathfrak B_\rho(h)-\mathfrak W_\rho(h)|
       \leq2\delta\mathcal N_\delta(h)Z_{2,\rho},
 \qquad
 Z_{2,\rho}=\sum_\omega\frac{m_\omega}{1+\gamma_\omega^2}
                  \ll_K\log(2+q_\rho),
\end{equation}
with $\mathcal N_\delta$ from Theorem~\ref{thm:band-bound}.
By \eqref{eq:artin-induced-trace}, this bounds the deviation of
the uncompressed trace limit from $\mathfrak W_\rho$.
\end{corollary}

\begin{proof}
The intersection $\ker\chi\cap j(\ker\chi)j^{-1}$ is open and
normal in $G_\Q$, so the induction factors through a finite group.
In the two coset lines its restriction to $G_K$ is
$\chi\oplus\chi^\sigma$, and $j$ exchanges these lines.
Since the characters differ, some element has distinct
eigenvalues on the two lines. Any invariant line must be one
of them, which contradicts the exchange. This proves
irreducibility and dimension two.

Artin induction gives \eqref{eq:artin-induced-L} prime by prime
\citep[Theorem~5.10 and Lemma~5.12, pp.~59--61]{horawaLfunctions}.
Explicitly, write $W_\rho$ for the representation space,
$I_p$ for inertia at a rational prime $p$, and
$\operatorname{Frob}_p$ for arithmetic Frobenius modulo inertia.
The local assertion is
\[
 \det(1-p^{-s}\rho(\operatorname{Frob}_p)\mid W_\rho^{I_p})^{-1}
   =\prod_{\mathfrak p\mid p}
                 (1-\chi(\mathfrak p)(N\mathfrak p)^{-s})^{-1}.
\]
Inertia invariants remove a factor exactly when the inducing
character is ramified there. An unramified character at the
prime above three still contributes its degree-one factor;
ramification of $K/\Q$ does not delete it.
The conductor induction formula
\citep[Proposition~5.58(2), p.~78]{horawaLfunctions} gives
$\operatorname{cond}(\rho)=|D_K|N\mathfrak f_\chi=q_\rho$,
where $D_K=-3$ is the field discriminant.
The nonprincipal Hecke function is entire.

The two eigenvalues of $j$ are $+1,-1$.
Gamma duplication consequently gives the standard completion
\[
 q_\rho^{s/2}\Gamma_\R(s)\Gamma_\R(s+1)L_\Q(s,\rho)
                =2c_\rho^{-s}\Gamma(s)L_K(s,\chi).
\]
The constant two does not change the functional equation or
zero divisor, and $c_\rho$ is exactly the Hankel parameter.

On ideals $\chi^\sigma(\mathfrak a)=\chi(\sigma\mathfrak a)$.
Conjugation preserves norms and conjugates conductors, so
reindexing the theta sum gives $E_{\chi^\sigma}^K=E_\chi^K$.
Their raw selections and full Gram projections are identical.
Thus one representative of the orbit $\{\chi,\chi^\sigma\}$
represents the Artin factor; summing both sectors counts it twice.

The normalized idele norms satisfy
$|N_{K/\Q}x|_\Q=|x|_K$. The inherited variable is therefore the
absolute norm $r$, with $t=\log r$ and Haar measure $dr/r$.
There is no rescaling of $t$ by two: the single norm channel
has leading term $2\log\Lambda$, not $4\log\Lambda$.

The zeros in \eqref{eq:artin-induced-L} coincide with full
multiplicity, and the defined projection is precisely the
Hecke projection of Theorem~\ref{nh:theorem}. That theorem
proves the trace and measure assertions and the polynomial
mass bound. Bessel's inequality also gives the pointwise
positivity in \eqref{eq:artin-induced-measures}.
There are no principal polar coordinates because $\chi\ne1$.
The unitary functional equation and coefficient conjugation
reflect a zero $\omega$ to $1-\overline\omega$; this does not
require $\chi^\sigma=\overline\chi$.
The half-plane thus confines the divisor to width $\delta$.
Theorem~\ref{thm:band-bound} and Corollary~\ref{cor:hecke-band}
give \eqref{eq:artin-induced-band}.
\end{proof}

This is a degree-two Galois representation, not a nonabelian
character of the abelian group $C_\Q$. Its cutoff is inherited
from the $K$-Hecke norm channel; no $\rho$-isotypic component
of a $C_\Q$ representation is constructed.
The cubic Kummer example above is an instance: its $S_3$
permutation representation is $1\oplus\rho$, so its irreducible
Artin factor has degree two despite the cubic field.
This direct norm-variable identification is a subcase of
Corollary~\ref{hnew:odd-artin}, which also covers the odd
two-dimensional representations not induced from this field.

Even two-dimensional representations require a different
archimedean transform. For an arbitrary irreducible degree-three representation,
an index-two character induction cannot apply. Even a monomial
induction from a cubic field requires a cutoff theorem for that
different field and its archimedean factors.
In general a zeros-only trace also requires exclusion of poles
in the relevant strip; Brauer induction has possibly negative
integer coefficients and yields a quotient of Hecke functions.
Even given holomorphy and an automorphic theta realization,
the Gram control, normalized zero-jet tails, endpoint
closures, graph approximation, local limit, polynomial frequency
mass, and trace transfer still have to be proved for that cutoff.
If the windowed construction of Section~\ref{glnw:section} is used,
one must additionally prove the normalization comparison
$\Tr((\widetilde Q-S_LQS_L)(V(h)\otimes I_d))\to0$.
Exact finite-section positivity alone supplies none of these
limits for those remaining representation classes.
